\chapter{Methodology}
\label{chap:method}

\draftnote{Every quantity in this chapter is taken from a stored manifest, validation record, or experiment output rather than from prose. Values that appear only in a written report and not in a machine-readable artefact are attributed to that report explicitly. The independent review described in Section~\ref{sec:method-review} has since been completed and the benchmark frozen; results against the frozen labels are reported in Chapter~\ref{chap:results}.}

\noindent
This chapter specifies the system the study measures, in the order in which a
question passes through it. Figure~\ref{fig:pipeline} shows that path in full;
the sections that follow take it one component at a time, and each records what
was fixed on development data and what was left open.

\begin{figure}[p]
  \centering
  \resizebox{\textwidth}{!}{%% Figure: end-to-end pipeline architecture.
%% Requires tikz with libraries arrows.meta, positioning, calc.
%% Colours aivancityNavy / aivancityGold / aivancityLight / aivancityGray /
%% draftRed are defined in main.tex.
\begin{tikzpicture}[
  font=\footnotesize,
  every node/.append style={execute at begin node={%
    \hyphenpenalty=10000\exhyphenpenalty=10000\relax}},
  stage/.style={
    draw=aivancityNavy, line width=0.6pt, rounded corners=2pt,
    fill=white, align=center, inner sep=4pt, text width=8.4cm
  },
  input/.style={
    stage, fill=aivancityLight
  },
  outcome/.style={
    draw=aivancityGold, line width=0.9pt, rounded corners=2pt,
    fill=white, align=center, inner sep=4pt, text width=2.62cm
  },
  aside/.style={
    draw=aivancityGray, line width=0.5pt, dashed, rounded corners=2pt,
    fill=white, align=left, inner sep=4pt, text width=3.2cm,
    text=aivancityGray, font=\footnotesize
  },
  flow/.style={-{Latex[length=4pt,width=4pt]}, draw=aivancityNavy, line width=0.7pt},
  bus/.style={draw=aivancityNavy, line width=0.7pt},
  side/.style={-{Latex[length=4pt,width=4pt]}, draw=aivancityGray, line width=0.5pt, dashed},
]

%% --------------------------------------------------------------- inputs ---
\node[input, text width=7.6cm] (corpus) at (0,0)
  {\textbf{Source corpus}\\[1pt]
   Reg.\ (EU) No 376/2014, with its implementing and delegated acts and the Guidance Material\\[1pt]
   $\rightarrow$ 109 legal parents, 1{,}535 retrieval children};

\node[input, text width=3.9cm, anchor=west] (question) at ($(corpus.east)+(5mm,0)$)
  {\textbf{Question}\\[1pt] frozen benchmark, 98 records};

\node[stage, text width=7.6cm, below=7mm of corpus] (index)
  {\textbf{Index}\\[1pt] child \texttt{retrieval\_text}, the hierarchy-enriched form};

%% ------------------------------------------------------------ retrieval ---
\coordinate (rowA) at ($(index.south)+(0,-11mm)$);
\node[stage, text width=3.9cm, anchor=north east] (dense) at ($(rowA)+(-3mm,0)$)
  {\textbf{Dense retrieval}\\[1pt] hosted embeddings, cosine};
\node[stage, text width=3.9cm, anchor=north west] (bm25) at ($(rowA)+(3mm,0)$)
  {\textbf{Lexical retrieval}\\[1pt] BM25, $k_1 = 1.5$, $b = 0.75$};

\node[stage, anchor=north] (rrf) at ($(dense.south)!0.5!(bm25.south)+(0,-7mm)$)
  {\textbf{Rank fusion} --- equal-weight RRF ($c\,{=}\,60$), no tuned parameter\\[1pt]
   $\rightarrow$ top-20 candidate pool};

\node[stage, anchor=north] (rerank) at ($(rrf.south)+(0,-7mm)$)
  {\textbf{Cross-encoder reranking} --- local MiniLM, no hosted call\\[1pt]
   reorders the pool; recall at depth 20 is unchanged by construction};

%% -------------------------------------------------------------- context ---
\node[stage, anchor=north] (context) at ($(rerank.south)+(0,-7mm)$)
  {\textbf{Context assembly} --- no gold label is consulted\\[2pt]
   \emph{compact}: 5 children, $\leq 1{,}200$ words \quad$\cdot$\quad
   \emph{enriched}: same 5 anchors $+$ siblings at radius 1, same budget\\[2pt]
   13-field block whitelist; a block carries the child's \texttt{text}, never its \texttt{retrieval\_text}};

%% ----------------------------------------------------------- generation ---
\node[stage, anchor=north] (gen) at ($(context.south)+(0,-7mm)$)
  {\textbf{Evidence-constrained generation}\\[2pt]
   closed JSON schema --- claims, supports, abstention flag and reason --- with a
   twelve-rule prompt contract; the model writes source ids, never citations};

%% --------------------------------------------------------- verification ---
\node[stage, anchor=north] (verify) at ($(gen.south)+(0,-7mm)$)
  {\textbf{Structural verification} --- deterministic, no gold label\\[2pt]
   schema validity \quad$\cdot$\quad every source id present in the assembled context
   \quad$\cdot$\quad every quote contiguous in the evidence it cites};

%% ------------------------------------------------------------- outcomes ---
\coordinate (rowO) at ($(verify.south)+(0,-11mm)$);
\node[outcome, anchor=north] (clarify) at (rowO)
  {\textbf{Clarify or retry}\\[1pt] verification failed, readings compete};
\node[outcome, anchor=north east] (accept) at ($(rowO)+(-2.85cm,0)$)
  {\textbf{Answer released}\\[1pt] claims with resolved citations};
\node[outcome, anchor=north west] (abstain) at ($(rowO)+(2.85cm,0)$)
  {\textbf{Abstain}\\[1pt] referral to a human reviewer};

%% ---------------------------------------------------------------- asides ---
\node[aside, anchor=west] (resolve) at ($(gen.east)+(6mm,0)$)
  {The application, not the model, resolves a source identifier to its canonical citation.};
\node[aside, anchor=east] (eval) at ($(context.west)+(-6mm,0)$)
  {Gold labels are read only when the run is scored, never inside the loop.};

%% ----------------------------------------------------------------- edges ---
\draw[flow] (corpus) -- (index);

\coordinate (busL) at ($(dense.north)+(0,5mm)$);
\coordinate (busR) at ($(bm25.north)+(0,5mm)$);
\coordinate (busM) at ($(busL)!0.5!(busR)$);
\draw[bus] (busL) -- (busR);
\draw[bus] (index.south) -- (busM);
\draw[bus] (question.south) |- (busM);
\draw[flow] (busL) -- (dense.north);
\draw[flow] (busR) -- (bm25.north);

\draw[flow] (dense.south) -- ++(0,-3.5mm) -| (rrf.north);
\draw[flow] (bm25.south) -- ++(0,-3.5mm) -| (rrf.north);
\draw[flow] (rrf) -- (rerank);
\draw[flow] (rerank) -- (context);
\draw[flow] (context) -- (gen);
\draw[flow] (gen) -- (verify);
\draw[flow] (verify.south) -- ++(0,-5.5mm) -| (accept.north);
\draw[flow] (verify.south) -- ++(0,-5.5mm) -| (clarify.north);
\draw[flow] (verify.south) -- ++(0,-5.5mm) -| (abstain.north);
\draw[side] (resolve.west) -- (gen.east);
\draw[side] (eval.east) -- (context.west);

\end{tikzpicture}
}
  \caption[The evaluated pipeline, end to end]{The evaluated pipeline, end to
  end. Solid arrows carry the runtime path; dashed boxes record two properties
  that hold throughout. The generator is given source identifiers and returns
  them, so a citation string is produced by the application and never by the
  model; and no gold label is read at any point between the question arriving
  and an outcome being chosen. Retrieval parameters are those fixed on the
  development partition (Section~\ref{sec:method-retrieval}).}
  \label{fig:pipeline}
\end{figure}

\section{Research Design}
\label{sec:method-design}

\subsection{Design Principles}

The study is a staged ablation. Each stage changes one major component while holding the corpus, the benchmark identifiers, and the source-grouped partition fixed, so that a change in measured performance can be attributed to the component that moved. Four principles govern the design and are enforced in code rather than by convention.

\begin{enumerate}
  \item \textbf{Retrieval is evaluated before generation.} A failed answer can originate in missing evidence, poor ranking, context construction, generation, or verification. Aggregating these into one score destroys the information needed to choose an intervention, which is the object of study here.
  \item \textbf{One component changes at a time.} Where two changes are of interest, they are run as separate experiments with separate manifests rather than as a combined configuration.
  \item \textbf{Selection happens on development data only.} Model choice, fusion constants, and any threshold are fixed on the development partition. The test partition is used for held-out reporting and never for tuning.
  \item \textbf{Failed and rejected configurations are retained.} Negative ablations remain in the repository and in the decision log rather than being silently replaced, so that the reported configuration can be read against the alternatives that were tried.
\end{enumerate}

Listing~\ref{lst:split-guard} shows the third principle as it exists in the
code. The generation driver defaults to the development partition and refuses
the held-out one outright unless a flag is passed whose name is a statement the
author has to be willing to make. A convention can be forgotten under deadline;
this cannot.

\begin{lstlisting}[style=thesis,caption={The held-out split cannot be generated by accident. Source: \texttt{scripts/run\_generation.py}, lines 181--186.},label={lst:split-guard}]
    if args.split in ("test", "all") and not args.i_have_frozen_calibration:
        raise SystemExit(
            f"Refusing to generate on '{args.split}': the held-out split is only "
            "scored once, after calibration is frozen on development data. "
            "Re-run with --i-have-frozen-calibration if that is genuinely done."
        )
\end{lstlisting}

\subsection{Experiment Ladder}

The ladder proceeds from the frozen historical baseline to the interventions under study. Each rung isolates a single decision.

The stages are set out in Table~\ref{tab:ladder}.

\begin{table}[htbp]
  \centering
  \caption{Experiment ladder. Each stage alters one component relative to the stage above it.}
  \label{tab:ladder}
  \small
  \begin{tabularx}{\textwidth}{@{}l l l X@{}}
    \toprule
    \textbf{ID} & \textbf{Chunking} & \textbf{Retrieval} & \textbf{Component under test} \\
    \midrule
    E0 & Flat legacy & Dense (3-small) & Frozen historical baseline \\
    E1 & Parent--child & Dense (3-small) & Effect of hierarchy-aware chunking \\
    E2 & Parent--child & Dense (3-large) & Effect of embedding capacity \\
    E3 & Parent--child & BM25 + dense, \ac{RRF} & Effect of a second retrieval channel \\
    E4 & Parent--child & E3 + cross-encoder & Effect of reranking a fixed pool \\
    E5 & Best configuration & E4 & Context assembly, generation, verification \\
    E6 & Best configuration & E4 & Optional constrained multi-step retrieval \\
    \bottomrule
  \end{tabularx}
\end{table}

Operationally the sequence is: audit the historical corpus and evaluation artefacts; construct a deterministic parent--child corpus; map historical questions to exact legal evidence; compare two dense embedding models on identical inputs; add local BM25 and equal-weight \ac{RRF}; rerank the fused top-20 pool with a local cross-encoder; assemble bounded evidence contexts; generate structured answers and verify claims, quotes, and citations; and calibrate clarification, retrieval expansion, abstention, and human-review triggers.

The experiment registry was reconciled with the stored manifests on 25 July 2026. It records E1--E4 as completed candidate runs, resolves E3 to \texttt{text-embedding-3-large} under D014, and distinguishes completed E5 context assembly from the generation phase, which was still blocked at that date and has since been run (Chapter~\ref{chap:results}).

\subsection{Unit of Analysis}

The unit of analysis is the question. For each question the study records the retrieved candidate set at several depths, the reranked ordering, the assembled context, the generated structured answer, the outcome of each verification check, and the escalation level at which the question was resolved. Aggregate figures are computed over questions within a partition; no figure is reported without its denominator.

\section{Corpus Construction}
\label{sec:method-corpus}

\subsection{Defects in the Historical Corpus}

The historical corpus contained 132 comparatively flat chunks derived from the same regulatory source. Three defect classes were identified in the audit and had to be repaired before any retrieval error could be interpreted as a model limitation.

First, chunk size was uncontrolled: some chunks exceeded 30{,}000 characters while the generator received only the first 1{,}500 characters of the selected chunk. A source could therefore be counted as retrieved while the supporting paragraph was absent from the generation context, which makes retrieval success and answer success only loosely coupled. Second, document identity was not preserved: the Easy Access Rules publication consolidates several instruments, and similarly numbered provisions from different instruments could collide. Third, legal type was not reliably separated, so binding regulation text and \ac{GM} were at risk of being treated as equivalent evidence.

The third defect is the most consequential for this application. Guidance material explains how a requirement may be met; it does not itself impose the obligation. A system that cites guidance where the regulation governs produces an answer that is superficially reasonable and legally wrong.

\subsection{Source Snapshot and Provenance}

Three source artefacts are recorded in \texttt{data/manifests/source\_manifest.csv} with their SHA-256 digests and marked \texttt{immutable\_historical}: the \ac{EASA} Easy Access Rules XML export used as the primary corpus, the consolidated regulation text as the authoritative reference, and the Easy Access Rules PDF as a rendering reference for manual checks~\cite{eu2014reg376,easa2022easyaccess376}. Only the XML export is parsed; the other two support manual verification.

\draftnote{The manifest records the primary corpus with \texttt{version = 2022-12}, while the file name, the build manifest, and the corpus audit all describe it as a 27 September 2023 export. Both are correct --- it is the December 2022 revision of the Easy Access Rules, exported in September 2023 --- but the chapter should state both explicitly rather than choose one.}

The build manifest for the promoted corpus records the source path, the source digest, the schema version \texttt{parent\_child\_v1}, the experiment identifier, start and completion timestamps, the validation summary, and the digests of the emitted parent and child files. Downstream experiment manifests inherit those digests, so any change to the corpus invalidates the lineage of every result computed from it.

\subsection{Legal Identity}

The parser preserves the identity of each underlying instrument rather than treating the consolidated publication as one document. The promoted corpus spans five document identities: the base regulation, three amending or implementing instruments, and the guidance material issued for the base regulation. Each legal parent records its document, regulation, type, and version; its formal title and legal-unit type; its hierarchy path, parent title, and canonical citation; its cross-references and source range; and the ordered identifiers of its retrievable children.

This identity layer is what prevents an article numbered~6 in one instrument from being confused with an article numbered~6 in another, and it is what keeps regulation and guidance distinguishable at retrieval time rather than only at reading time.

\subsection{Parent--Child Representation}

A parent is a coherent legal unit: a complete article, an annex section, a recitals block, or a guidance section. A child is an atomic retrievable unit: a paragraph, list item, table row, table header, or callout. The design follows the finding that finer indexing granularity improves retrieval, particularly for rare entities~\cite{chen2024densex}, while retaining the structural context that makes a fragment of legal text interpretable.

Table~\ref{tab:schema-fields} lists the fields each unit stores and which of them are optional.

\begin{table}[htbp]
  \centering
  \caption{Stored fields. Parents carry fourteen fields, of which only cross-references are optional; all eleven child fields are required.}
  \label{tab:schema-fields}
  \small
  \begin{tabularx}{\textwidth}{@{}l X@{}}
    \toprule
    \textbf{Record} & \textbf{Fields} \\
    \midrule
    Parent &
      \texttt{parent\_id}, \texttt{document\_id}, \texttt{document\_version}, \texttt{regulation\_id},
      \texttt{document\_type}, \texttt{hierarchy\_path}, \texttt{canonical\_citation}, \texttt{title},
      \texttt{text}, \texttt{source\_location}, \texttt{cross\_references} (optional),
      \texttt{char\_count}, \texttt{word\_count}, \texttt{child\_ids} \\
    \addlinespace
    Child &
      \texttt{child\_id}, \texttt{parent\_id}, \texttt{sequence}, \texttt{canonical\_citation},
      \texttt{unit\_type}, \texttt{text}, \texttt{retrieval\_text}, \texttt{char\_count},
      \texttt{word\_count}, \texttt{marker\_path}, \texttt{source\_location} \\
    \bottomrule
  \end{tabularx}
\end{table}

Two field distinctions carry methodological weight. The child stores both \texttt{text} and \texttt{retrieval\_text}: the former is the evidence exactly as it appears in the source and is what a verifier and a human reader see, while the latter is a metadata-enriched form used for indexing. Separating them means that hierarchy can improve retrieval without contaminating the evidence that a quote is checked against. Table rows repeat the governing table and column headers in \texttt{retrieval\_text} for the same reason: the row becomes findable without the displayed evidence being altered.

Short legal units are not merged to satisfy a generic minimum length. An atomic item remains interpretable through its stored hierarchy, and merging would reintroduce exactly the boundary confounding that the repair was intended to remove (decision D007).

\subsection{Deterministic Validation}

Validation is implemented as fifteen distinct checks --- eleven critical and four warning --- applied at corpus, parent, child, and cross-reference level. Critical checks reject an empty corpus or empty child set, duplicate parent or child identifiers, empty parent or child text, empty retrieval text, a missing regulation identifier, a parent with no children, an orphaned child, an article whose title does not begin with ``Article'', and any mismatch between a parent's declared child ordering and the children that actually reference it. Warning checks flag an unusual guidance title, a parent above 40{,}000 characters, a child above 4{,}000 characters, and a table row that carries no column structure while exceeding twenty words.

The child-order check is the one that most directly protects downstream work: because context assembly expands to neighbouring siblings by position, an ordering mismatch would silently return the wrong neighbours.

Table~\ref{tab:e1-corpus} reports the composition of the corpus that passed validation.

\begin{table}[htbp]
  \centering
  \caption{Validated E1 corpus composition, from the stored validation record.}
  \label{tab:e1-corpus}
  \rowcolors{2}{aivancityLight}{white}
  \begin{tabular}{lr}
    \toprule
    \textbf{Unit} & \textbf{Count} \\
    \midrule
    Legal parents & 109 \\
    Retrievable children & 1,535 \\
    \midrule
    ARTICLE parents & 36 \\
    GM parents & 59 \\
    ANNEX parents & 10 \\
    RECITALS parents & 4 \\
    \midrule
    List-item children & 823 \\
    Paragraph children & 477 \\
    Callout children & 122 \\
    Table-row children & 105 \\
    Table-header children & 8 \\
    \bottomrule
  \end{tabular}
\end{table}

The promoted corpus reports zero critical issues and zero warnings. Child length ranges from 2 to 1{,}615 characters with a mean of 192.61. The corpus audit additionally reports 87 units below 20 characters and 8 units above 1{,}000 characters; these two counts appear in the audit report only and are not recorded in the validation output. A second independent build produced identical parent, child, and validation bytes, and the builder refuses to overwrite a non-empty output directory unless forced.

\section{Benchmark Construction}
\label{sec:method-benchmark}

\subsection{Legacy Lineage}

The historical dataset contains 300 synthetic question--answer records generated over the 132 flat chunks. Because the old source field records a chunk rather than a legal unit, it cannot be treated as exact evidence. Every record was therefore mapped into the E1 hierarchy under a four-stage precedence: normalised text equality; legacy text wholly contained in a new parent; a single new parent wholly contained in a legacy chunk; and, last, fuzzy token coverage with cosine ranking. Crucially, the fourth stage never accepts a mapping --- it marks the record as requiring review.

Table~\ref{tab:legacy-lineage} reports the outcome of that mapping over the 300 historical records.

\begin{table}[htbp]
  \centering
  \caption{Legacy parent-lineage audit over 300 historical records.}
  \label{tab:legacy-lineage}
  \begin{tabular}{lr}
    \toprule
    \textbf{Disposition} & \textbf{Records} \\
    \midrule
    Automatically mapped & 291 \\
    Manually accepted & 4 \\
    Excluded as invalid & 4 \\
    Rewrite required & 1 \\
    \midrule
    Total & 300 \\
    \bottomrule
  \end{tabular}
\end{table}

At chunk level, 113 of the 132 legacy chunks mapped automatically and 19 required review; of those 19, only 4 were used by question records lacking an automatic mapping, and the 9 affected records were resolved individually and recorded with their resolutions. This asymmetry matters: the number of ambiguous chunks overstates the number of ambiguous evaluation records by a factor of roughly two.

\subsection{Evidence Audit of the Shared Evaluation Set}

The 100 question identifiers shared by the stored few-shot, \ac{RAG}, and \ac{LoRA} outputs were then linked to exact child evidence. This audit was \ac{AI}-assisted, and the stored decision counts record exactly what was done to each record: 78 retained without text change, 21 rewritten and evidence-linked, and 1 excluded as invalid. The retained candidate benchmark contains 99 records, each identifying a document-qualified parent and at least one selected evidence child.

The excluded record, \texttt{QA\_0003}, was generated from concatenated table-of-contents text; its reference answer merely refers the reader to a section number and is not evidence-grounded. All three baseline systems failed its citation criterion, so its removal does not favour any system. Because the historical results were computed over 100 questions, both views are reported: the immutable $N=100$ baseline for continuity, and the $N=99$ sensitivity view for any claim that depends on evidence-linked gold data (decision D009).

\subsection{Independent Multi-Reviewer Validation}
\label{sec:method-review}

The audit described above cannot certify itself. The same \ac{AI}-assisted process that proposed corrections is not in a position to confirm them, and the wider literature gives direct grounds for caution: model-generated evaluation data requires human anchoring before its scores can be trusted~\cite{saadfalcon2024ares}, and model judges carry documented position, verbosity, and self-enhancement biases~\cite{zheng2023judging}. A separate validity gate is therefore mandatory before any figure computed on this benchmark is treated as final.

The review package contains 33 records selected deterministically: all 22 corrected or excluded audit exceptions; all three hybrid top-20 residuals, two of which already appear among the exceptions; and ten unchanged controls stratified by source class and split at two per stratum. The sample spans 16 ARTICLE, 13 GM, and 4 ANNEX records, and 16 development, 16 test, and 1 excluded record. The required identifiers are pinned in a requirements manifest together with the digests of the audited benchmark, the child corpus, and the split manifest, so that the review cannot be applied to a different version of the data than the one it was performed on.

For each record a reviewer records question clarity, answer support, citation correctness, evidence completeness, corrections, exclusions, and alternative acceptable evidence sets. Multi-child alternatives are represented as logical evidence routes rather than as a single union that must always be retrieved together, so that a system retrieving a legally sufficient alternative is not scored as having failed (decision D021).

The review is conducted by a panel of practising aviation professionals rather than by the author alone (decisions D023 and D024). Three properties of the protocol are fixed before any response is collected. First, a record is auto-decided only when at least two reviewers have rated it. Second, the aggregation rule is a strict majority on the overall decision and on each of the four validity checks, with \texttt{accept} requiring all four checks to resolve affirmatively. Third, a majority verdict of ``accept with correction'' is never applied automatically: the panel identifies the problem and the author supplies the corrected fields, because a corrected evidence set cannot be produced by a multiple-choice instrument.

Agreement is reported rather than asserted, using chance-corrected statistics: Cohen's coefficient for two raters~\cite{cohen1960coefficient}, Fleiss's generalisation for three or more~\cite{fleiss1971measuring}, with magnitudes interpreted against the conventional benchmarks and their arbitrariness acknowledged~\cite{landis1977measurement}. Records on which qualified reviewers disagree are reported individually rather than resolved by majority and absorbed into the aggregate. Disagreement between practitioners about whether a provision supports an answer is evidence that the question is ambiguous, and treating it as noise would discard the signal that identifies precisely the cases the escalation policy exists to handle.

The freeze command validates the exact 33 required identifiers, reviewer role and date, allowed decision values, evidence identifiers, source digests, and correction consistency. It then applies corrections, preserves provenance, excludes rejected records transparently, regenerates the parent-grouped splits, and writes SHA-256 manifests. It refuses to run against an incomplete review and refuses to overwrite a non-empty output directory.

\draftnote{This gate has since been passed. The review closed at the pre-registered cutoff, the 19 contested records were adjudicated by the author, and the benchmark was frozen at 98 records; every retrieval figure in Chapter~\ref{chap:results} is computed against those frozen labels.}

\subsection{Parent-Grouped Partitioning}

Random splitting at question level would leak near-identical evidence whenever several questions originate from one article or guidance section, and such leakage inflates apparent generalisation. The benchmark is therefore partitioned by legal parent, following the practice of reporting evidence provenance alongside answers~\cite{petroni2021kilt}.

The partitioner is deterministic without a random number generator. Groups are ordered by descending size, with ties broken by the SHA-256 digest of a string seed concatenated with the group identifier, and are then assigned greedily to whichever partition brings the running development count closest to a target fraction of records. Because the seed is a string consumed by a hash rather than an RNG state, the split reproduces exactly across platforms and library versions. A separate assertion raises an error if any group's questions carry more than one partition label.

The resulting partition is given in Table~\ref{tab:split}.

\begin{table}[htbp]
  \centering
  \caption{Candidate parent-grouped evaluation split, from the stored split summary.}
  \label{tab:split}
  \begin{tabular}{lrr}
    \toprule
    \textbf{Partition} & \textbf{Questions} & \textbf{Parent groups} \\
    \midrule
    Development & 42 & 12 \\
    Test & 57 & 48 \\
    \midrule
    Total & 99 & 60 \\
    \bottomrule
  \end{tabular}
\end{table}

Two features of Table~\ref{tab:split} require comment, since both look like errors and are not. The requested development fraction is 0.4, which for 99 records implies a target of 40; the greedy rule lands on 42 because groups are assigned whole and the target cannot generally be hit exactly. The development partition holds 42 questions across only 12 parents while the test partition holds 57 across 48, because the largest groups are ordered first and are therefore assigned first. Recorded group leakage is zero.

The split and the review sample use deterministic string seeds:

\begin{itemize}
  \item \texttt{aviation-rag-legacy100-v1} for the parent-grouped partition;
  \item \texttt{aviation-rag-independent-review-v1} for the review sample.
\end{itemize}

\noindent Both are consumed by SHA-256 assignment; no integer random-number-generator seed is used for these operations.

\section{Retrieval Architecture}
\label{sec:method-retrieval}

\subsection{Dense Retrieval}

Two hosted embedding models process identical inputs --- the 1{,}535 hierarchy-enriched child retrieval texts and the 99 candidate questions. Each model-specific run consumed 172{,}821 input tokens, for 345{,}642 across both. Vectors are cached with their input digests, model names, and request metadata, so that a change to a label does not force re-embedding when the underlying text is unchanged. Cosine similarity ranks child vectors for each query.

The larger model was selected on the development partition only, before the test comparison was opened (decision D014). No in-domain fine-tuning of the embedding model was performed; the design relies on the finding that general-purpose contrastively trained retrievers transfer to unseen domains without task-specific training~\cite{izacard2022contriever}.

\subsection{Lexical Retrieval}

BM25 is implemented locally with $k_1=1.5$ and $b=0.75$. For query $Q$ and child document $D$,

\begin{equation}
\operatorname{BM25}(D,Q) =
\sum_{q \in Q}
\operatorname{IDF}(q)\,
\frac{f(q,D)(k_1+1)}
     {f(q,D)+k_1\left(1-b+b\frac{|D|}{\operatorname{avgdl}}\right)},
\label{eq:bm25}
\end{equation}

where $f(q,D)$ is the term frequency, $|D|$ is the document length in tokens, and $\operatorname{avgdl}$ is the corpus mean. Inverse document frequency uses the non-negative form $\operatorname{IDF}(q) = \log\!\left(1 + \frac{N - n_q + 0.5}{n_q + 0.5}\right)$, and each term's contribution is additionally scaled by its query frequency; the formulation follows the standard treatment of term weighting and inverted indexing~\cite{manning2008ir}. Ties are broken by corpus position, so rankings are deterministic. The channel defined by Equation~\ref{eq:bm25} supplies exact-terminology matching for article references, acronyms, deadlines, and repeated regulatory expressions~\cite{robertson2009bm25}.

Tokenisation is deliberately minimal: text is lower-cased and split on the pattern \texttt{[a-z0-9]+}, with no stopword removal, stemming, or lemmatisation. Two consequences must be stated plainly rather than discovered by a reader. A citation such as ``376/2014'' tokenises into two separate numeric tokens, and a subparagraph reference such as ``4(1)(a)'' into four; the lexical channel therefore matches the components of a legal reference rather than the reference as a unit. And because the character class is restricted to unaccented ASCII after lower-casing, accented characters act as token separators, which bounds what this implementation could do on French-language queries without modification. Learned sparse retrieval would address the first limitation by design~\cite{formal2021splade}; it is recorded here as future work rather than as an implemented alternative.

\subsection{Rank Fusion}

The selected dense ranking and the BM25 ranking are combined by equal-weight \ac{RRF}, using Equation~\ref{eq:rrf} with $c = 60$. Ranks are one-based, duplicates within a single input ranking are ignored, and ties are resolved first by the best rank achieved in any input and then lexicographically by identifier, so that the fused ordering is reproducible across platforms. The constant was fixed before held-out reporting and was not tuned on test data.

Two alternatives were evaluated on development data and rejected: unequal fusion weights, and boosting candidates by parent mass. Neither improved the top-20 candidate objective, and parent boosting consistently reduced recall. Both are retained as documented negative ablations (decision D017).

The choice of \ac{RRF} over a tuned score-level combination is a deliberate trade of performance for methodological safety. A convex combination of normalised scores has been shown to outperform \ac{RRF} both in and out of domain, and rank-based fusion necessarily discards score magnitude~\cite{bruch2024fusion}. A fusion function with no tuned parameters, however, cannot leak information from the held-out partition, which in a study whose central claim concerns measurement validity is worth more than the marginal recall forgone. The trade-off is revisited in the limitations.

\subsection{Cross-Encoder Reranking}

The fused top-20 candidates are reranked by a cross-encoder that scores each query--passage pair jointly, following the established result that direct query--passage interaction substantially improves ranking over independently encoded representations~\cite{nogueira2019reranking}. The model is a distilled English passage-ranking cross-encoder~\cite{wang2020minilm} trained on the MS MARCO passage collection~\cite{nguyen2016msmarco}, run with a maximum sequence length of 512 and a batch size of 64 on a local CUDA device. The run scored 1{,}980 query--passage pairs over 99 queries, and the manifest records the model identifier, the exact model revision commit, the framework version, the device, and the digest of the cached score set. No question or regulatory passage was sent to a hosted inference endpoint during reranking.

The training distribution of this reranker must be named rather than assumed away. MS MARCO consists of short, factoid-oriented web search queries; the questions in this benchmark are longer, structurally explicit, and drawn from a specialist regulatory sublanguage. Any improvement is therefore obtained despite a domain shift, and any residual failure has to be read with that shift in view. A multilingual reranker was considered and deferred on scope and model-size grounds (decision D018).

The reranking implementation preserves each candidate's original hybrid rank alongside its new score, so that recovery and regression relative to the pre-reranking ordering can be measured per question rather than only in aggregate. A reranker can repair ordering only when the relevant evidence is already present in the candidate pool; recall at depth 20 is therefore expected to be unchanged, and that invariance is what distinguishes a ranking failure from a candidate-generation failure.

\section{Context Construction and the Generation Contract}
\label{sec:method-context}

\subsection{Context Selection}

Two deterministic context variants are assembled from the reranked ordering, neither of which consults gold labels at any point. The compact variant selects the top five children under a 1{,}200-word budget with no neighbour expansion. The enriched variant selects the same five anchors, expands to adjacent sibling units within a radius of one, and applies the same 1{,}200-word budget. An earlier build of the variant, over the 99 pre-review questions, used a 1{,}600-word budget instead; the difference makes no selection difference on this corpus, for the reason given below. Budgets are measured in stored child word counts.

Three behaviours of the selection procedure should be understood precisely, because each would otherwise look like a defect. The first anchor is admitted unconditionally, even if it alone exceeds the budget, so that no question is ever left with an empty context. A candidate skipped for budget does not stop the scan, so a lower-ranked but shorter child may still be admitted --- selection is best-fit by rank, not a prefix of the ranking. And neighbours receive no budget exemption and are drawn strictly from siblings within the same parent, never across parents, so expansion cannot silently introduce a different legal unit.

Bounding the context is not a resource optimisation. Models use evidence unreliably when it sits in the middle of a long context, and performance degrades as context length grows even in models designed for long inputs~\cite{liu2024lostmiddle}. A tight, ordered evidence budget is therefore a correctness measure, and it is what makes reranking worth its cost.

Each context block exposes a fixed whitelist of thirteen fields: source identifier, parent identifier, role, anchor identifier, retrieval rank, sequence, canonical citation, document identifier, document type, document version, hierarchy path, unit type, evidence text, and word count. Notably the block carries the child's \texttt{text}, not its \texttt{retrieval\_text}; the metadata-enriched indexing form is never shown to the generator, so a quote can be checked against exactly what the model was given. Parent title and text, source locations, and marker paths are deliberately withheld.

Table~\ref{tab:context-stats} reports what the selection procedure actually produces on this corpus.

\begin{table}[htbp]
  \centering
  \caption{Observed context statistics for the enriched variant, as built over the 98 frozen questions --- the artefact that generation reads. The earlier build over the 99 pre-review questions gave 6.36 neighbours and 473.3 evidence words per question under the same selection rule.}
  \label{tab:context-stats}
  \begin{tabular}{lr}
    \toprule
    \textbf{Quantity} & \textbf{Value} \\
    \midrule
    Anchors per question (mean) & 5.00 \\
    Neighbours per question (mean) & 6.27 \\
    Sources per question (mean / median / max) & 11.27 / 11 / 15 \\
    Evidence words per question (mean / median / max) & 463.8 / 451.5 / 1,090 \\
    Questions hitting the word budget & 0 \\
    \bottomrule
  \end{tabular}
\end{table}

The maximum of fifteen sources is the structural ceiling for five anchors at radius one. No candidate was skipped for budget, and the largest observed context is 1{,}090 words against a 1{,}200-word ceiling. Because nothing was skipped, a larger ceiling could not have admitted any further unit: the enriched variant is bounded by structure rather than by the budget in this corpus, which is why the earlier 1{,}600-word setting and the present one select the same evidence.

\subsection{Answer Schema}

The generator returns a single JSON object validated against a closed schema that forbids additional properties. Five top-level properties are required: the answer text; an array of claims; an abstention flag; an abstention reason, which is either null or one of five controlled values (ambiguous question, conflicting sources, insufficient evidence, outside corpus, verification failed); and a clarification question, which may be null.

Each claim carries a non-empty identifier, its text, and at least one support. Each support carries a non-empty source identifier and a non-empty verbatim quote. The minimum-item constraint on supports is the structural core of the contract: a claim with no support cannot validate.

Two limits of the schema are recorded here rather than left for a reader to find. It does not encode the cross-field rule that the abstention reason must be null when the abstention flag is false --- that rule is stated in the prompt and must be checked separately. And it does not restrict source identifiers to the supplied set; that restriction is applied at runtime against the identifiers actually present in the assembled context.

\subsection{Prompt Contract}

The prompt instructs the model to answer only from the supplied source blocks, not to rely on unstated legal knowledge, and not to invent article numbers, citations, obligations, exceptions, deadlines, or consequences. Twelve numbered rules operationalise this: decompose the answer into the smallest material claims; give each claim a unique identifier; support every claim with one or more supplied source identifiers; attach a short verbatim quote to each support; never write an unsupplied source identifier; do not generate display citations at all; abstain with the appropriate controlled reason when evidence is insufficient, when the question admits several plausible legal meanings, when supplied sources materially conflict, or when the question falls outside the corpus; set the abstention reason to null when not abstaining; and distinguish binding regulation text from guidance where the distinction affects the answer.

The sixth rule carries the most weight. The application, not the model, resolves source identifiers to canonical citations, so a model cannot produce a valid-looking citation string by itself. This follows the finding that retrieve-then-read pipelines outperform post-hoc citation attachment on both correctness and attribution~\cite{bohnet2022attributed}, and it directly targets the failure mode measured in deployed systems, where citations are frequently present but do not support the statements they are attached to~\cite{liu2023verifiability}.

\section{Verification and Selective Answering}
\label{sec:method-verify}

\subsection{Structural Verification}

Structural verification is deterministic and applies three checks. Schema validation rejects malformed output. The allowed-source check rejects any source identifier not present in the assembled context. The quote check requires each supporting quote to occur in the evidence text of the source it cites, compared after collapsing whitespace, trimming, and case folding. Punctuation and hyphenation are not normalised, so the check is strict on exactly the features that distinguish one legal formulation from another.

The outcome is a decision among three states: structurally accepted, abstained, or returned for retry or abstention. One asymmetry in the implementation should be understood by anyone reading the results. A claim can be marked unsupported --- because its source identifier is unknown or its quote is blank --- without contributing an entry to the quote-error list. The decision state can therefore be ``structurally accepted'' while a specific claim is not among the supported claims. Support coverage must accordingly be read from the list of supported claim identifiers rather than inferred from the decision.

Listing~\ref{lst:contiguity} is that quote check. Normalisation removes
typography and punctuation, which the corpus inherits from a PDF extraction and
which carry no legal meaning; the containment test is then applied to the
normalised word sequence, so a quote is accepted only where its words appear
\emph{contiguously} in the evidence. Contiguity is not relaxed, and it is the
condition that rejects a quote assembled from separated fragments.

\begin{lstlisting}[style=thesis,caption={The check the argument of this thesis rests on. Source: \texttt{src/aviation\_rag/generation/verification.py}, lines 72--80; reproduced with the byte-level rule beside it in Appendix~\ref{app:code}.},label={lst:contiguity}]
def _normalize_for_matching(value: str) -> str:
    """Typography- and punctuation-insensitive; word order is preserved."""
    value = NON_WORD.sub(" ", _normalize_typography(value))
    return WHITESPACE.sub(" ", value).strip().casefold()


def quote_is_grounded(quote: str, evidence_text: str) -> bool:
    """Does the quote's word sequence appear contiguously in the evidence?"""
    return _normalize_for_matching(quote) in _normalize_for_matching(evidence_text)
\end{lstlisting}

\subsection{The Limits of Structural Verification}

Structural checks establish provenance, not entailment. A quote may be reproduced perfectly and still fail to support the claim built on it, and material omission --- an answer that states an obligation while omitting the condition that suspends it --- is invisible to a containment test. This is the extrinsic case that dominates the risk profile in generation over authoritative text~\cite{ji2023hallucination}, and the empirical record in deployed legal systems shows that retrieval grounding reduces but does not eliminate it~\cite{magesh2025hallucination}.

Accordingly, every non-abstaining answer is flagged for separate semantic review, and claim support and completeness are evaluated independently of the structural pass rate. The formal criterion for that review is attributability: a claim counts as supported only if a competent reader would affirm that, according to the cited source, the claim holds~\cite{rashkin2023attribution}. Reporting a structural pass rate as though it were a groundedness rate would repeat precisely the conflation this thesis sets out to measure.

\subsection{Selective Answering}

The system records accepted answers, retries, clarification requests, and abstentions, and selective performance is reported as a risk--coverage relationship rather than as a single accuracy figure~\cite{geifman2017selective}. The policy prefers a correct refusal to an unsupported regulatory answer, and this preference is a design commitment rather than an empirical finding.

The signal driving abstention is deliberately not the generator's own confidence. Model confidence is poorly calibrated under domain shift, and a decision function separate from the answering model outperforms thresholding the model's own score~\cite{kamath2020selective}. The triggers used here are therefore external and observable: retrieval score margin, disagreement between the lexical and dense channels, absence of any candidate above threshold, and the outcome of structural verification.

\section{Escalation Policy and Its Calibration}
\label{sec:method-escalation}

The escalation policy is the object of study rather than a component of the system, and it is specified as a ladder of levels with an explicit transition rule between them. Table~\ref{tab:escalation-method} sets out the levels and the condition for moving beyond each of them. The same ladder is redrawn in Figure~\ref{fig:escalation-policy}, in Chapter~\ref{chap:results}, once the results have fixed the evidential status of each transition.

\begin{table}[htbp]
  \centering
  \caption{Escalation levels and their trigger conditions. All thresholds are calibrated on development data only.}
  \label{tab:escalation-method}
  \small
  \begin{tabularx}{\textwidth}{@{}c l X@{}}
    \toprule
    \textbf{Level} & \textbf{Action} & \textbf{Trigger for moving beyond this level} \\
    \midrule
    0 & Dense retrieval only & Baseline; superseded by construction \\
    1 & Hierarchy-aware retrieval & Default after corpus repair \\
    2 & Hybrid candidates with fusion & Low fused score, or lexical/dense disagreement \\
    3 & Cross-encoder reranking & Small margin between top candidates \\
    4 & Evidence-constrained generation & Applied to all answered questions \\
    5 & Clarification or targeted revision & Verification failure, or several plausible readings \\
    6 & Abstention and human referral & Conflicting sources, no supporting evidence, unresolved ambiguity \\
    \bottomrule
  \end{tabularx}
\end{table}

Three constraints govern calibration. First, every trigger must be computable at runtime: gold-label agreement is unavailable in deployment and cannot appear in a trigger, however convenient it would be for reporting. Second, thresholds are fitted on the development partition and applied unchanged to the test partition. Third, the policy is evaluated not only on final accuracy but on whether it escalated the right questions, through escalation precision --- the proportion of escalated questions that would in fact have failed at the lower level --- and escalation recall --- the proportion of failing questions that were escalated. A policy that escalates everything achieves perfect recall and is worthless; reporting both quantities is what makes the claim of an economical policy falsifiable.

\section{Evaluation Measures}
\label{sec:method-measures}

\subsection{Retrieval Measures}

For a question $j$ with a set of acceptable evidence routes $\mathcal{G}_j$ and retrieved top-$k$ set $R_{j,k}$, evidence-route recall is

\begin{equation}
\operatorname{EvidenceRecall}_{j}@k =
\max_{G \in \mathcal{G}_j}
\frac{|R_{j,k} \cap G|}{|G|}.
\label{eq:evidence-recall}
\end{equation}

The maximum over routes in Equation~\ref{eq:evidence-recall} is what allows a legally sufficient alternative to count as success. Any-gold-child recall records whether at least one child belonging to any acceptable route appears; parent recall records whether a child of an acceptable legal parent appears. \ac{MRR} is

\begin{equation}
\operatorname{MRR}@k =
\frac{1}{N}\sum_{j=1}^{N}\frac{1}{r_j},
\label{eq:mrr}
\end{equation}

In Equation~\ref{eq:mrr}, $r_j$ is the rank of the first acceptable child, contributing zero when none occurs within $k$. \ac{nDCG} additionally rewards acceptable children appearing earlier, crediting graded rather than binary relevance and discounting by position~\cite{jarvelin2002cumulated}.

Each experiment stores these measures at cutoffs $k \in \{1, 3, 5, 10, 20\}$, separately for the development partition, the test partition, and all questions, together with the denominator $n$ in each view.

The experiment registry now uses the machine-readable keys written by the runs: any-gold-child recall, mean evidence recall, parent recall, \ac{MRR}, \ac{nDCG}, and the denominator $n$, at cutoffs 1, 3, 5, 10, and 20.

\subsection{Generation Measures}

Generation is evaluated through three separate families, which are never combined into one score.

Structural measures record the valid-schema rate, the allowed-source-only rate, the quote-verification rate, and canonical citation accuracy. Answer-quality measures record claim support, completeness, and exact legal-answer accuracy. Selective measures record coverage, risk at coverage, and appropriate-abstention rate. Operational measures --- latency, input and output tokens, and estimated cost --- are recorded alongside, since an escalation policy that improves accuracy at unbounded cost has not solved the stated problem.

Surface-overlap metrics are reported only as continuity with the historical baseline and are never presented as evidence of correctness. Recall-oriented n-gram overlap~\cite{lin2004rouge} correlates only weakly with human faithfulness judgements, while entailment-based measures correlate considerably better~\cite{maynez2020faithfulness}. Citation precision and recall are reported separately from answer correctness, following the benchmark practice established for citation-bearing generation~\cite{gao2023alce}.

\subsection{Statistical Treatment}

Proportions are reported with score-based confidence intervals rather than normal-approximation intervals, which are unreliable at the per-condition sample sizes here and at proportions near zero or one~\cite{wilson1927probable}. Comparisons between systems on the same questions are treated as paired, and the count of questions each system uniquely passes is reported alongside the difference in aggregate rates, because two systems can post similar totals while failing on disjoint sets. Subgroup results with small denominators are reported with their denominators visible and are not used to support general claims.

\section{Threats to Validity}
\label{sec:method-threats}

Table~\ref{tab:threats} states the principal threats to the validity of this study and the control applied to each.

\begin{table}[htbp]
  \centering
  \caption{Principal threats to validity and the control applied to each.}
  \label{tab:threats}
  \small
  \begin{tabularx}{\textwidth}{@{}l X@{}}
    \toprule
    \textbf{Threat} & \textbf{Control} \\
    \midrule
    Benchmark provenance & Questions and reference answers were machine-generated; a stratified sample is audited by a panel of practitioners and agreement is reported~\cite{saadfalcon2024ares}. \\
    Evidence leakage across splits & Partitioning by legal parent, with a programmatic assertion that no group spans both partitions. \\
    Contamination of public benchmarks & A benchmark built over one specific instrument rather than reused from a public test set~\cite{sainz2023contamination}. \\
    Reranker domain shift & The training distribution is stated, and reranking gains are interpreted as obtained despite that shift~\cite{nguyen2016msmarco}. \\
    Threshold tuning on test data & All selection and calibration performed on development; fusion uses an untuned constant. \\
    Metric conflation & Retrieval, structural, semantic, selective, and operational measures reported separately. \\
    Corpus defects mistaken for model limits & Corpus repaired and validated before any retrieval error is attributed to a model. \\
    Single-instrument scope & Results are stated for one regulation; no claim of generalisation to other instruments is made. \\
    Author involvement in review & Panel ratings reported with and without the author's own; contested records adjudicated openly and reported separately. \\
    \bottomrule
  \end{tabularx}
\end{table}

\section{Reproducibility and Audit Trail}
\label{sec:method-repro}

Experiment directories are immutable by default. Each manifest records input and output digests, parameters, software and framework versions, model identifiers and revisions, hardware, timestamps, and a status field. Dense vectors and cross-encoder scores are cached under their input digests so that a repeated run is verifiably the same computation rather than a similar one. Determinism is achieved through explicit tie-breaking and hash-derived ordering rather than through seeded random number generators, so that results reproduce across platforms and library versions.

The test strategy requires tests before data or results may be promoted, covering corpus construction, retrieval behaviour, evaluation semantics, and safety properties. The corpus audit records twelve parser and integrity tests passing at promotion, together with a byte-identical independent rebuild.

The corpus builders, the retrieval and generation drivers, the evaluator, the verification layer and the test suite are published in full at \url{https://github.com/MoSaleh47/Where-RAG-Breaks---Citation-Grounded-Question-Answering-and-Escalation-over-EU-Aviation-Regulations}, with a README describing the pipeline and the order in which the scripts are run; Appendix~\ref{app:code} reproduces the excerpts the argument depends on most directly.

Decisions are append-only and each reversal references the decision it supersedes. Rejected ablations remain documented. Every run computed before the freeze carries an explicit review-pending status in its manifest and is kept apart from the runs computed on frozen labels; the historical baseline is likewise preserved separately from the repaired pipeline, so that none of the three is silently conflated with another.

\section{Ethical and Professional Considerations}
\label{sec:method-ethics}

Three commitments constrain the method. Transmission of corpus and question text to a hosted embedding service was treated as requiring explicit authorisation, was refused until that authorisation was given, and is recorded with the token counts actually processed (decisions D012 and D013). Reranking runs locally precisely so that no project text reaches a hosted inference endpoint at that stage. And the \ac{AI}-assisted curation of the benchmark is disclosed in the reporting configuration and in Appendix~\ref{app:ai-disclosure} rather than presented as human-authored gold data.

The system is a decision-support research prototype. It is not a compliance decision-maker, and any output it produces requires verification by a qualified person before it is relied upon. This constraint is stated in the user-facing output and is tested as a safety property, not left as a disclaimer in the documentation.
